\documentclass[11pt]{article}
\usepackage[a4paper,margin=1in]{geometry}
\usepackage{fontspec}
\usepackage[boldfont=false]{xeCJK}
\setCJKmainfont{FandolSong-Regular.otf}
\usepackage{amsmath}
\usepackage{amssymb}
\usepackage{array}
\usepackage{booktabs}
\usepackage{graphicx}
\usepackage{adjustbox}
\usepackage{enumitem}
\setlist[itemize]{topsep=2pt,partopsep=0pt,parsep=0pt,itemsep=2pt,leftmargin=1.4em}
\setlist[enumerate]{topsep=2pt,partopsep=0pt,parsep=0pt,itemsep=2pt,leftmargin=1.6em}
\usepackage{parskip}
\usepackage{fvextra}
\fvset{breaklines=true,breakanywhere=true,fontsize=\small}
\setlength{\parindent}{0pt}

\begin{document}

\begin{center}
  {\LARGE\bfseries Human resource management}\\[0.35em]
  {\large A Level Business}
\end{center}
\vspace{0.6em}

\section*{What human resource management does}

\textbf{Human resource management} 人力资源管理 (HRM) is the part of a business that looks after its people. Good HRM makes sure the firm has enough staff, with the right skills, who are well trained and motivated.

The main roles of HRM are:

\begin{itemize}
\item \textbf{workforce planning} 人力规划 — working out how many staff, and what skills, the business will need in the future.

\item \textbf{recruitment} 招聘 and \textbf{selection} 甄选 — finding people to apply for jobs, then choosing the best.

\item \textbf{training} 培训 and \textbf{development} 培养 — improving staff skills now and helping them grow for the future.

\end{itemize}

\section*{Workforce planning, recruitment and selection}

Workforce planning looks at the firm's plans, then at the staff it already has, to find the gap. The gap may be filled by hiring, training, or moving staff.

To recruit, the firm first writes:

\begin{itemize}
\item a \textbf{job description} 职位描述 — the title, duties and tasks of the job.

\item a \textbf{person specification} 人员规格 — the skills, \textbf{qualifications} 资格 and qualities the right person should have.

\end{itemize}

The firm can recruit from inside or outside:

\begin{itemize}
\item \textbf{internal recruitment} 内部招聘 — filling the job with someone who already works there. It is cheaper and faster, and the person is known, but no new ideas come in.

\item \textbf{external recruitment} 外部招聘 — hiring someone from outside. It brings new skills and ideas, but costs more and is slower.

\end{itemize}

Selection methods include reading application forms and CVs, \textbf{interviews} 面试, tests, and group tasks.

\section*{Training and development}

Training improves the skills staff need to do their jobs well. Common types are:

\begin{itemize}
\item \textbf{induction training} 入职培训 — given to new staff so they learn how the firm works.

\item \textbf{on-the-job training} 在岗培训 — learning while doing the real job.

\item \textbf{off-the-job training} 脱产培训 — learning away from the workplace, such as a course.

\end{itemize}

Training costs time and money, but it raises quality, \textbf{productivity} 生产率 and motivation, and helps keep staff.

\section*{Ending employment}

Sometimes a worker has to leave.

\begin{itemize}
\item \textbf{redundancy} 裁员 happens when the job itself is no longer needed — for example, when a machine replaces it. It is not the worker's fault.

\item \textbf{dismissal} 解雇 happens when a worker is made to leave because of poor work or bad behaviour.

\end{itemize}

\section*{Rights and managing change}

Both sides have rights and duties. The \textbf{employer} 雇主 must follow the law on pay, safety, working hours and unfair dismissal. The \textbf{employee} 雇员 must work as agreed and follow safe, fair rules. These \textbf{rights} 权利 protect both sides.

Businesses change often — new technology, new markets, or \textbf{restructuring} 重组. The \textbf{management of change} 变革管理 in HR means planning the change, explaining it clearly, training staff for new roles, and listening to worries, so staff accept the change instead of resisting it.

\section*{Measuring the workforce}

Two simple measures help managers judge how the workforce is doing.

\textbf{Labour turnover} 员工流失率 shows the percentage of staff who leave in a year:

\[
\text{labour turnover} = \frac{\text{number of staff leaving}}{\text{average number employed}} \times 100\%
\]

High labour turnover can mean low pay, poor motivation or weak management. It raises the cost of recruiting and training.

\textbf{Labour productivity} 劳动生产率 shows the output made by each worker:

\[
\text{labour productivity} = \frac{\text{total output}}{\text{number of employees}}
\]

Higher labour productivity lowers the cost of each unit, so the firm can compete better.

\section*{Why motivation matters}

\textbf{Motivation} 激励 is the desire to work hard and do a good job. Motivated staff produce more and better work, take fewer days off, stay longer, and give better service. So motivation lowers costs and raises quality.

\section*{Theories of motivation}

Several thinkers tried to explain what motivates people. You should know the main idea of each.

\begin{center}
\adjustbox{max width=\linewidth}{%
\begin{tabular}{ll}
\toprule
\textbf{Thinker} & \textbf{Main idea} \\
\midrule
\textbf{Taylor} & \textbf{scientific management} 科学管理 — workers are mainly motivated by money. Pay them per unit made (\textbf{piece rate} 计件工资) and they work harder. \\
\textbf{Mayo} & \textbf{human relations} 人际关系 — workers are also motivated by attention, teamwork and feeling part of a group. \\
\textbf{Maslow} & \textbf{hierarchy of needs} 需求层次 — people meet needs in order, from low to high (see below). \\
\textbf{Herzberg} & \textbf{two-factor theory} 双因素理论 — some things cause unhappiness if missing, others truly motivate (see below). \\
\textbf{McClelland} & people are driven by three needs in different amounts: \textbf{achievement} 成就, power and affiliation (belonging). \\
\textbf{Vroom} & \textbf{expectancy theory} 期望理论 — people work hard only if they believe the effort will lead to a reward they want. \\
\bottomrule
\end{tabular}}
\end{center}

\subsection*{Maslow's hierarchy of needs}

Maslow said people try to meet lower needs first, then move up:

\begin{enumerate}
\item \textbf{physiological needs} 生理需求 — food, water, rest (met by pay).

\item safety needs — a safe job and workplace.

\item social needs — friendship and being part of a team.

\item \textbf{esteem needs} 尊重需求 — respect and \textbf{recognition} 认可.

\item \textbf{self-actualisation} 自我实现 — reaching your full potential.

\end{enumerate}

\subsection*{Herzberg's two factors}

\begin{itemize}
\item \textbf{hygiene factors} 保健因素 — things like pay, conditions and rules. If they are poor, staff are unhappy; but fixing them does not truly motivate.

\item \textbf{motivators} 激励因素 — things like achievement, responsibility and recognition. These give real, lasting motivation.

\end{itemize}

\section*{Financial motivators}

\textbf{Financial motivators} 经济激励 are rewards paid in money:

\begin{itemize}
\item \textbf{wages} 工资 — pay by the hour or week, often for \textbf{manual work} 体力劳动.

\item \textbf{salary} 薪水 — a fixed yearly amount, paid monthly.

\item piece rate — pay for each unit made.

\item \textbf{commission} 佣金 — pay based on how much you sell.

\item \textbf{bonus} 奖金 — an extra payment for good results.

\item \textbf{profit sharing} 利润分享 — staff get a share of the firm's profit.

\item \textbf{performance-related pay} 绩效工资 — extra pay for those who meet targets.

\item \textbf{fringe benefits} 附加福利 — non-cash extras, such as a company car or health care.

\end{itemize}

\section*{Non-financial motivators}

\textbf{Non-financial motivators} 非经济激励 reward staff without extra money, by making work more interesting and giving more control:

\begin{itemize}
\item \textbf{job enrichment} 工作丰富化 — giving more challenging tasks and responsibility.

\item \textbf{job rotation} 工作轮换 — moving between different tasks to reduce boredom.

\item \textbf{job enlargement} 工作扩大化 — adding more tasks at the same level.

\item \textbf{empowerment} 授权 — letting staff make their own decisions about their work.

\item \textbf{teamworking} 团队合作 — organising staff into teams that share goals.

\end{itemize}

\section*{Management and its functions}

\textbf{Management} 管理 means getting work done through other people to reach the firm's objectives.

\textbf{Fayol} said managers carry out five functions:

\begin{center}
\adjustbox{max width=\linewidth}{%
\begin{tabular}{ll}
\toprule
\textbf{Function} & \textbf{What the manager does} \\
\midrule
\textbf{planning} 计划 & set objectives and decide how to reach them \\
\textbf{organising} 组织 & arrange people and resources \\
\textbf{commanding} 指挥 & give clear instructions and guidance \\
\textbf{coordinating} 协调 & make sure all parts work together \\
\textbf{controlling} 控制 & check results against the plan and correct problems \\
\bottomrule
\end{tabular}}
\end{center}

\textbf{Mintzberg} said managers also play ten \textbf{roles} 角色 in three groups: interpersonal (working with people), informational (sharing information), and decisional (making decisions). This shows that real management is busy and varied, not just the five neat functions.

\section*{Leadership versus management}

\textbf{Leadership} 领导 and management are not the same.

\begin{itemize}
\item a manager plans, organises and controls the day-to-day work.

\item a leader sets a \textbf{vision} 愿景, inspires people, and guides change.

\end{itemize}

The best managers are also good leaders. Effective management matters because it raises productivity, keeps good staff, helps the firm adapt to change, and turns plans into real results.


\end{document}
